\section{Language Models and RNNs}
\label{sec:language-models}
\slidesrc{Language models \& RNNs: 1--101}

Until now we have studied language mostly by \emph{analysing} it: splitting it into
words, tagging them, parsing sentences, representing meanings. This section takes a
different point of view, the one behind every modern NLP system and every chatbot:
we build a \textbf{model that can predict language}. Given the beginning of a
sentence, which word comes next? Given a whole sentence, how likely is it?

A model that answers these questions is called a \textbf{language model} (LM). GPT and
the other Large Language Models are language models; they are just very large and
built with a more advanced architecture (the transformer). To understand them we
start from the beginning: what a language model is, the simplest way to build one by
counting (n-grams), how neural networks improve on counting, and how
\textbf{recurrent neural networks} (RNNs) read text of any length, together with the
problems that eventually led to their replacement.

No previous knowledge of neural networks is assumed: every tool (probability chain
rule, softmax, cross-entropy, gradient descent, backpropagation) is introduced when it
is first needed.

This section covers:
\begin{enumerate}[nosep]
  \item what a language model is and why it is useful
        (Section~\ref{sec:what-is-a-language-model});
  \item classical language models built by counting: n-grams, their sparsity and
        storage problems, and text generation (Section~\ref{sec:ngram-models});
  \item neural language models with a fixed window, word embeddings, the softmax and
        the cross-entropy loss (Section~\ref{sec:neural-language-models});
  \item recurrent neural networks: architecture, training, backpropagation through
        time, text generation and evaluation with perplexity
        (Section~\ref{sec:recurrent-neural-networks});
  \item the problems of RNNs: vanishing and exploding gradients, gradient clipping and
        the ideas that solve them (Section~\ref{sec:vanishing-exploding-gradients}).
\end{enumerate}

\subsection{What Is a Language Model?}
\label{sec:what-is-a-language-model}
\slidesrc{Language models \& RNNs: 2--11}

\subsubsection{Predicting and evaluating language}
\slidesrc{Language models \& RNNs: 3--6}

A \textbf{language model} is a mathematical model that describes human language. It
should let us do two things:
\begin{itemize}[nosep]
  \item \textbf{predict} language: continue a text that has been started;
  \item \textbf{evaluate} language: say how ``natural'' or likely a piece of text is.
\end{itemize}

How can we build such a model? Two facts make it hard to write it down as a fixed set of
rules:
\begin{itemize}[nosep]
  \item \textbf{language is ambiguous}: the same beginning can be continued in many
        reasonable ways (``she went to \ldots'' the shop? work? bed?);
  \item \textbf{language is evolving}: new words and new ways of saying things appear
        all the time, so any fixed list of rules goes out of date.
\end{itemize}
Both facts push in the same direction: instead of saying \emph{which} continuation is
correct, we say \emph{how likely} each continuation is. We therefore use
\textbf{probabilistic models of language}.

\subsubsection{Next-word prediction}
\slidesrc{Language models \& RNNs: 7--9}

The fundamental goal of a language model is \textbf{prediction}, and the simplest form
of prediction is guessing the \textbf{next word}. We number the words of a sentence by
their position: $\vect x^{(1)}$ is the first word, $\vect x^{(2)}$ the second, and so on
(the superscript $(t)$ is a position, ``time step $t$'', not a power).

\begin{center}
\begin{tikzpicture}[word/.style={font=\small, anchor=base, inner sep=2pt},
                    lab/.style={font=\scriptsize, nbgray}]
  \foreach \w/\x [count=\i] in {Anna/0, was/1.25, tired/2.45, so/3.55, she/4.5,
                                went/5.5, to/6.55} {
    \node[word] (w\i) at (\x,0) {\w};
    \node[lab] at (\x,-0.5) {$\vect x^{(\i)}$};
  }
  \node[word, draw=nbred, dashed, rounded corners, minimum width=1.1cm] (w8) at (7.85,0) {\strut ?};
  \node[lab, nbred] at (7.85,-0.55) {$\vect x^{(8)}$};
\end{tikzpicture}
\captionof{figure}{Next-word prediction: having read the first seven words, the model
must say which word is likely to come in position 8.}
\end{center}

Given the words $\vect x^{(1)}, \vect x^{(2)}, \ldots, \vect x^{(t)}$ read so far, the
model computes the \textbf{probability distribution of the next word}
$\vect x^{(t+1)}$:
\[
  p\bigl(\vect x^{(t+1)} \mid \vect x^{(t)}, \ldots, \vect x^{(1)}\bigr).
\]

\notinslides
A few words on this notation, since it is used everywhere in this section.
\begin{itemize}[nosep]
  \item The \textbf{vocabulary} $V$ is the set of all words the model knows, and $|V|$
        is its size (tens of thousands of words for a real model).
  \item A \textbf{probability distribution over the next word} is a list of $|V|$
        numbers, one per word of the vocabulary, each between 0 and 1, which add up to
        1. For example, after ``she went to'' a model could give \emph{the}: 0.5,
        \emph{work}: 0.25, \emph{school}: 0.2, \emph{bed}: 0.05, and 0 to all other
        words.
  \item The vertical bar ``$\mid$'' reads ``given'':
        $p(A \mid B)$ is the probability of $A$ \emph{knowing that} $B$ happened.
        What is written after the bar is called the \textbf{context} or the
        \textbf{conditioning}.
\end{itemize}

Language models are part of daily life: the slide shows a search engine completing a
query while it is being typed (each suggestion is a likely continuation of the
characters typed so far) and a chat assistant writing an answer word by word.

\subsubsection{Evaluating a text: the chain rule}
\slidesrc{Language models \& RNNs: 10}

An equivalent goal is to \textbf{evaluate} a text, i.e.\ assign it a probability. Given
the words $\vect x^{(1)}, \ldots, \vect x^{(T)}$ ($T$ is the length of the text), we want
$p(\vect x^{(1)}, \ldots, \vect x^{(T)})$: the probability that a text is exactly this
sequence of words. Why is this ``equivalent'' to next-word prediction? Because of the
\textbf{chain rule of probability}.

The starting point is the definition of conditional probability:
\[
  p(A, B) = p(A \mid B)\, p(B),
\]
``the probability that both $A$ and $B$ happen is the probability of $B$, times the
probability of $A$ once we know $B$ happened''. Applying it again and again we can peel
a sentence one word at a time, from the end. For three words:
\begin{align*}
  p(\vect x^{(1)}, \vect x^{(2)}, \vect x^{(3)})
  &= p(\vect x^{(3)} \mid \vect x^{(2)}, \vect x^{(1)})\; p(\vect x^{(1)}, \vect x^{(2)})
     \reason{$A = \vect x^{(3)}$, $B = (\vect x^{(1)}, \vect x^{(2)})$} \\
  &= p(\vect x^{(3)} \mid \vect x^{(2)}, \vect x^{(1)})\; p(\vect x^{(2)} \mid \vect x^{(1)})\;
     p(\vect x^{(1)}) \reason{again, on $p(\vect x^{(1)}, \vect x^{(2)})$}
\end{align*}
Reading the last line from right to left: probability of the first word, times the
probability of the second word given the first, times the probability of the third
given the first two. In general:

\begin{nbfact}{Chain rule for a text}
\begin{align*}
  p(\vect x^{(1)}, \ldots, \vect x^{(T)})
  &= p(\vect x^{(1)}) \times p(\vect x^{(2)} \mid \vect x^{(1)}) \times \cdots \times
     p(\vect x^{(T)} \mid \vect x^{(T-1)}, \ldots, \vect x^{(1)}) \\
  &= \prod_{t=1}^{T} \underbrace{p(\vect x^{(t)} \mid \vect x^{(t-1)}, \ldots,
     \vect x^{(1)})}_{\text{next-word prediction!}}
\end{align*}
($\prod$ is the product symbol, the multiplicative analogue of $\sum$; for $t = 1$ the
context is empty and the factor is just $p(\vect x^{(1)})$.)
\end{nbfact}

So the two goals are really one: a model that can predict the next word can evaluate
any text (multiply the next-word probabilities of all its words), and a model that can
evaluate texts can predict the next word (compare the probabilities of the possible
continuations). This identity is exact: no approximation has been made yet.

\begin{nbexample}{probability of a short sentence}
\notinslides
Suppose a model says $p(\textit{the}) = 0.06$ for the first word,
$p(\textit{cat} \mid \textit{the}) = 0.01$ and
$p(\textit{sleeps} \mid \textit{the cat}) = 0.05$. Then
\[
  p(\textit{the cat sleeps}) = 0.06 \times 0.01 \times 0.05 = 0.00003 .
\]
The number is tiny, as is the probability of any particular sentence (there are
astronomically many possible sentences), but it can be \emph{compared}: a good model
gives \emph{the cat sleeps} a much higher probability than \emph{sleeps cat the}.
\end{nbexample}

\subsubsection{Why care about language models?}
\slidesrc{Language models \& RNNs: 11}

\begin{nbdef}{Language model}
A \textbf{language model} is a probabilistic model of natural language: it assigns a
probability to every sequence of words, or equivalently gives, for every context, the
probability distribution of the next word.
\end{nbdef}

\begin{itemize}[nosep]
  \item Language modelling is a good \textbf{benchmark} for general NLP models: to
        predict the next word well a model must know grammar, meaning and facts about
        the world.
  \item Every task that needs to \textbf{estimate the probability of a text} can be
        built on top of an LM: speech recognition, handwriting recognition, spell
        checking, machine translation, summarisation, dialogue, \ldots
  \item Today \textbf{everything else in NLP is built on top of language modelling}:
        GPT is a language model.
\end{itemize}

\begin{nbintuition}[how a probability of text helps other tasks]
\notinslides
A speech recogniser hears a sound that could be ``recognise speech'' or ``wreck a nice
beach''. Both match the audio; the language model says the first is a far more likely
English sentence, so it wins. A spell checker that sees ``I went to the sea \emph{shor}''
proposes \emph{shore} rather than \emph{short} because $p(\textit{shore} \mid
\textit{the sea})$ is much higher. In both cases the LM provides the knowledge of
``what people usually say''.
\end{nbintuition}

\subsection{Classical Language Models: n-grams}
\label{sec:ngram-models}
\slidesrc{Language models \& RNNs: 12--34}

Before neural networks, language models were built by \textbf{counting}. This section
shows how, why it works surprisingly well, and why it eventually hits a wall.

\subsubsection{n-grams and the Markov assumption}
\slidesrc{Language models \& RNNs: 13--16}

\begin{nbdef}{n-gram}
An \textbf{n-gram} is a sequence of $n$ consecutive words of a text.
\end{nbdef}

For the sentence ``The quick brown fox jumps over the lazy dog'':
\begin{center}
\small
\begin{tabular}{@{}ll@{}}
\toprule
$n$ & n-grams \\
\midrule
1 (\textbf{unigrams})  & \emph{the}, \emph{quick}, \emph{brown}, \emph{fox}, \ldots \\
2 (\textbf{bigrams})   & \emph{the quick}, \emph{quick brown}, \emph{brown fox}, \emph{fox jumps}, \ldots \\
3 (\textbf{trigrams})  & \emph{the quick brown}, \emph{quick brown fox}, \emph{brown fox jumps}, \ldots \\
4 (\textbf{four-grams})& \emph{the quick brown fox}, \emph{quick brown fox jumps}, \ldots \\
\bottomrule
\end{tabular}
\end{center}

An \textbf{n-gram language model}:
\begin{enumerate}[nosep]
  \item collects the \textbf{frequencies} (counts) of all n-grams in a large text
        corpus;
  \item uses these frequencies to predict the next word.
\end{enumerate}

\paragraph{The Markov assumption.}
The chain rule (Section~\ref{sec:what-is-a-language-model}) needs the probability
of the next word given the
\emph{whole} history. That is impossible to estimate by counting: most long histories
appear only once in any corpus, or never. An n-gram model makes a simplifying
assumption, named after the mathematician Markov (we met the same idea in Hidden Markov
Models, Section~\ref{sec:hmm-pos-tagging}).

\begin{nbdef}{Markov assumption (n-gram model)}
The next word depends only on the $n-1$ words just before it:
\[
  p\bigl(\vect x^{(t+1)} \mid \vect x^{(t)}, \ldots, \vect x^{(1)}\bigr) \approx
  p\bigl(\vect x^{(t+1)} \mid
  \underbrace{\vect x^{(t)}, \ldots, \vect x^{(t-n+2)}}_{n-1 \text{ words}}\bigr).
\]
(Counting the context: from position $t-n+2$ to position $t$ there are
$t - (t-n+2) + 1 = n-1$ words. Together with the predicted word they form one n-gram.)
\end{nbdef}

For example, a \textbf{bigram} model ($n = 2$) looks only at the previous word. The
probability of ``the quick brown fox'' is decomposed as
\[
  p(\texttt{<s>}) \, p(\textit{the} \mid \texttt{<s>}) \, p(\textit{quick} \mid \textit{the})
  \, p(\textit{brown} \mid \textit{quick}) \, p(\textit{fox} \mid \textit{brown}) \,
  p(\texttt{</s>} \mid \textit{fox}).
\]
\texttt{<s>} and \texttt{</s>} are two special tokens marking the \textbf{beginning} and
the \textbf{end of the sentence}. They let the model learn which words tend to
\emph{start} a sentence ($p(\textit{the} \mid \texttt{<s>})$ is high) and which tend to
\emph{end} it. Every sentence starts with \texttt{<s>}, so $p(\texttt{<s>}) = 1$.

\paragraph{Estimating the probabilities by counting.}
By the definition of conditional probability, the probability we need is a ratio:
\[
  p\bigl(\vect x^{(t+1)} \mid \vect x^{(t)}, \ldots, \vect x^{(t-n+2)}\bigr) =
  \frac{p\bigl(\vect x^{(t+1)}, \vect x^{(t)}, \ldots, \vect x^{(t-n+2)}\bigr)}
       {p\bigl(\vect x^{(t)}, \ldots, \vect x^{(t-n+2)}\bigr)} .
\]
We do not know these probabilities, but we have a corpus, a finite sample of language.
The \textbf{finite-sample approximation} replaces each probability with how often the
sequence occurs in the corpus:
\[
  p\bigl(\vect x^{(t+1)} \mid \vect x^{(t)}, \ldots, \vect x^{(t-n+2)}\bigr) \approx
  \frac{\overbrace{\text{count}\bigl(\vect x^{(t+1)}, \vect x^{(t)}, \ldots,
        \vect x^{(t-n+2)}\bigr)}^{\text{an n-gram}}}
       {\underbrace{\text{count}\bigl(\vect x^{(t)}, \ldots,
        \vect x^{(t-n+2)}\bigr)}_{\text{an $(n-1)$-gram}}} .
\]
In words: \emph{among all the times the context appeared, in what fraction was it
followed by this word?} The quality of the approximation depends on the
\textbf{size of the corpus}: the more text, the closer the fractions get to the true
probabilities.

\begin{nbexample}{a four-gram model}
\slidesrc{Language models \& RNNs: 17--19}
We want to continue ``Anna was tired so she went to \underline{\quad}'' with a
\textbf{four-gram} model ($n = 4$). It looks only at the last $n - 1 = 3$ words,
\emph{she went to}, and discards the rest:
\[
  p(x \mid \textit{she went to}) \approx
  \frac{\text{count}(\textit{she went to } x)}{\text{count}(\textit{she went to})} .
\]
Suppose the corpus contains:
\begin{center}
\small
\begin{tabular}{@{}lrl@{}}
\toprule
sequence & count & estimated probability of $x$ \\
\midrule
\emph{she went to}        & 1000 & (the context) \\
\emph{she went to the}    & 500  & $500/1000 = 0.50$ \\
\emph{she went to work}   & 250  & $250/1000 = 0.25$ \\
\emph{she went to school} & 200  & $200/1000 = 0.20$ \\
\emph{she went to bed}    & 50   & $50/1000 = 0.05$ \\
\bottomrule
\end{tabular}
\end{center}
The four probabilities add up to 1 (here every occurrence of the context is followed by
one of these four words). The model predicts \emph{the} as most likely and \emph{bed} as
least likely.

\textbf{Was it a good idea to discard the longer context?} No: ``Anna was
\emph{tired}'' is exactly the clue that makes \emph{bed} the most sensible
continuation, but the model never saw it, because it fell outside the window of three
words.
\end{nbexample}

\begin{nbexample}{a bigram model from a toy corpus}
\notinslides
The corpus has three sentences:
\begin{center}
\texttt{<s>} I like cats \texttt{</s>} \qquad
\texttt{<s>} I like dogs \texttt{</s>} \qquad
\texttt{<s>} you like cats \texttt{</s>}
\end{center}
Counts we need: \texttt{<s>} appears 3 times, followed twice by \emph{I} and once by
\emph{you}; \emph{I} appears 2 times, always followed by \emph{like}; \emph{like}
appears 3 times, followed twice by \emph{cats} and once by \emph{dogs}; \emph{dogs}
appears once, followed by \texttt{</s>}. So
\[
  p(\textit{I} \mid \texttt{<s>}) = \tfrac23, \quad
  p(\textit{like} \mid \textit{I}) = \tfrac22 = 1, \quad
  p(\textit{dogs} \mid \textit{like}) = \tfrac13, \quad
  p(\texttt{</s>} \mid \textit{dogs}) = \tfrac11 = 1 .
\]
The probability of the sentence ``I like dogs'' is
\begin{align*}
  p(\textit{I like dogs})
  &= p(\textit{I} \mid \texttt{<s>})\, p(\textit{like} \mid \textit{I})\,
     p(\textit{dogs} \mid \textit{like})\, p(\texttt{</s>} \mid \textit{dogs}) \\
  &= \tfrac23 \cdot 1 \cdot \tfrac13 \cdot 1 = \tfrac29 \approx 0.22 .
\end{align*}
Two observations that lead to the next topic:
\begin{itemize}[nosep]
  \item ``you like dogs'' \emph{never} appears in the corpus, yet the model gives it
        $\frac13 \cdot 1 \cdot \frac13 \cdot 1 = \frac19$: by cutting sentences into
        short pieces, the model \textbf{generalises} to new sentences;
  \item ``I hate cats'' gets probability \textbf{0}, because the bigram \emph{I hate}
        was never seen. One unseen bigram is enough to kill the whole sentence.
\end{itemize}
\end{nbexample}

\subsubsection{Sparsity problems}
\slidesrc{Language models \& RNNs: 20--26}

\paragraph{Sparsity problem 1: unseen continuations.}
For many words $x$, the sequence ``she went to $x$'' \textbf{never occurs} in the
corpus (``she went to Reykjavik'', ``she went to dinner'' may simply be absent). Their
count is 0, so their probability is 0: the model says they are \emph{impossible}. Many
words with probability 0 make the LM \textbf{inflexible}.

\textbf{Mitigation: smoothing.} Add a small value $\epsilon$ to the count of every
word $x \in V$, so that nothing is ever exactly 0:
\[
  p(x \mid \textit{she went to}) =
  \frac{\text{count}(\textit{she went to } x) + \epsilon}
       {\text{count}(\textit{she went to}) + \epsilon |V|} .
\]

\begin{nbwarning}[the denominator must change too]
The slide adds $\epsilon$ only to the numerator. Then the probabilities no longer add
up to 1: summing the numerators over all $|V|$ words gives
$\text{count}(\textit{she went to}) + \epsilon|V|$, not
$\text{count}(\textit{she went to})$. Adding $\epsilon|V|$ to the denominator, as
above, fixes this. With $\epsilon = 1$ the method is called \textbf{add-one} or
\textbf{Laplace smoothing}.
\end{nbwarning}

\begin{nbexample}{add-one smoothing on the toy corpus}
\notinslides
In the toy corpus above, a word can be followed by one of the 6 tokens
\emph{I}, \emph{you}, \emph{like}, \emph{cats}, \emph{dogs}, \texttt{</s>}, so
$|V| = 6$; and $\text{count}(\textit{I}) = 2$. With $\epsilon = 1$:
\[
  p(\textit{like} \mid \textit{I}) = \frac{2 + 1}{2 + 6} = \frac38, \qquad
  p(\textit{cats} \mid \textit{I}) = \frac{0 + 1}{2 + 6} = \frac18 ,
\]
and the same $\frac18$ for each of the other 4 unseen tokens.
Check: $\frac38 + 5 \cdot \frac18 = 1$. \emph{like} is still by far the most likely, but
no word is impossible any more. The price: probability mass has been taken from the
words that were actually seen ($1 \to \frac38$) and spread over unseen ones.
\end{nbexample}

\paragraph{Sparsity problem 2: unseen contexts.}
What if ``she went to'' \textbf{never appeared} in the corpus? Then the denominator is
0, and we cannot compute the probability at all.

\textbf{Mitigation: backoff.} Condition on a shorter context, just ``went to'' (or
``she went'', as the slide writes), and if that is also missing, shorter still, down to
the unigram probability of $x$. Shorter contexts are more frequent, so they are more
likely to have been observed.

\begin{nbfact}{Increasing n makes sparsity worse}
A longer n-gram is a more specific sequence, so it is seen more rarely. With a larger
$n$ the model uses more context, but more of its counts are 0 and smoothing/backoff
are needed more often. In practice $n \le 5$.
\end{nbfact}

\paragraph{An example of sparsity.}
The \textbf{Berkeley Restaurant Project} dataset contains 9332 sentences of people
asking about restaurants, with a vocabulary of $|V| = 1446$ words. The table shows the
bigram counts for 8 of these words: the entry in row \emph{want}, column \emph{to} is
how many times \emph{want} was immediately followed by \emph{to}.

\begin{center}
\small
\begin{tabular}{@{}l*{8}{r}@{}}
\toprule
 & \emph{i} & \emph{want} & \emph{to} & \emph{eat} & \emph{chinese} & \emph{food} &
   \emph{lunch} & \emph{spend} \\
\midrule
\emph{i}       & 5  & 827 & 0  & 9   & 0 & 0 & 0 & 2   \\
\emph{want}    & 2  & 0  & 608 & 1   & 6  & 6  & 5  & 1   \\
\emph{to}      & 2  & 0  & 4   & 686 & 2  & 0 & 6  & 211 \\
\emph{eat}     & 0 & 0  & 2   & 0  & 16 & 2  & 42 & 0  \\
\emph{chinese} & 1  & 0  & 0  & 0  & 0 & 82 & 1  & 0  \\
\emph{food}    & 15 & 0  & 15  & 0  & 1  & 4  & 0 & 0  \\
\emph{lunch}   & 2  & 0  & 0  & 0  & 0 & 1  & 0 & 0  \\
\emph{spend}   & 1  & 0  & 1   & 0  & 0 & 0 & 0 & 0  \\
\bottomrule
\end{tabular}
\captionof{table}{Bigram counts in the Berkeley Restaurant Project (row = first word,
column = second word). Half of the 64 entries are 0, even among frequent words.}
\label{tab:berkeley-bigrams}
\end{center}

32 of the 64 cells are zero. Some zeros are real impossibilities (\emph{want want}),
but others are perfectly good English that simply did not occur (\emph{lunch food}?
\emph{spend lunch}?). And this is with \textbf{bigrams}: there are
$1446^2 \approx 2.1$ million possible bigrams, far more than a corpus of 9332 sentences
can show. With $n > 2$ it is much worse.

\subsubsection{Storage problems}
\slidesrc{Language models \& RNNs: 27}

An n-gram model must \textbf{store the counts} of all n-grams seen in the corpus. As $n$
grows, or as the corpus grows, this can become huge. With $|V| = 1446$ there are
\[
  1446^4 \approx 4.37 \times 10^{12}
\]
possible four-grams: naively storing a count for each of them, even at just one byte per
count, requires about \textbf{4 TB} of memory, for a toy vocabulary. Real vocabularies
have tens of thousands of words. (Storing only the n-grams actually seen saves a lot,
but on a large corpus that is still billions of entries.)

\subsubsection{Generating text with an n-gram model}
\slidesrc{Language models \& RNNs: 28--34}

A language model can also \textbf{generate} text: start from a context, compute the
distribution of the next word, \textbf{sample} a word from it, append it to the context
and repeat.

\begin{nbdef}{Sampling from a distribution}
\notinslides
\textbf{Sampling} a word means drawing it at random \emph{with the probabilities given
by the model}: a word with probability 0.3 is picked, on average, 3 times out of 10. It
is like rolling a die whose faces are the words of the vocabulary and whose faces have
different sizes. Sampling (instead of always taking the most likely word) gives varied
text rather than the same sentence every time.
\end{nbdef}

The slides use a \textbf{trigram} model built on the \textbf{Reuters corpus} (1.7
million words of financial news). Starting from ``today the'':

\begin{center}
\begin{tikzpicture}
\begin{axis}[xbar, width=6.5cm, height=4.2cm, bar width=7pt, xmin=0, xmax=0.2,
    symbolic y coords={emirate,italian,price,bank,company}, ytick=data,
    y tick label style={font=\small\itshape}, x tick label style={font=\scriptsize,
    /pgf/number format/fixed}, nodes near coords, nodes near coords style={font=\scriptsize},
    every node near coord/.append style={/pgf/number format/fixed,
    /pgf/number format/precision=3},
    title={\small $p(x \mid \textit{today the})$}, axis x line*=bottom, axis y line*=left,
    enlarge y limits=0.15]
  \addplot[fill=nbblue!60, draw=nbblue] coordinates
    {(0.153,company) (0.153,bank) (0.077,price) (0.039,italian) (0.039,emirate)};
\end{axis}
\end{tikzpicture}
\captionof{figure}{The five most likely words after ``today the'' in the trigram model
(the rest of the probability is spread over many other words).}
\end{center}

\begin{enumerate}[nosep]
  \item Context \emph{today the} $\to$ distribution: \emph{company} 0.153,
        \emph{bank} 0.153, \emph{price} 0.077, \emph{italian} 0.039,
        \emph{emirate} 0.039, \ldots\ Sample: \emph{price}.
  \item Text \emph{today the price}; context (last 2 words) \emph{the price} $\to$
        \emph{of} 0.308, \emph{for} 0.050, \emph{it} 0.046, \emph{to} 0.046,
        \emph{is} 0.031, \ldots\ Sample: \emph{of}.
  \item Text \emph{today the price of}; context \emph{price of} $\to$ \emph{the} 0.072,
        \emph{18} 0.043, \emph{oil} 0.043, \emph{its} 0.036, \emph{gold} 0.018, \ldots
        and so on.
\end{enumerate}
Note that at each step only the last $n - 1 = 2$ words are used as context: the
beginning \emph{today} is forgotten as soon as it leaves the window.

\paragraph{The effect of $n$.}
The slide shows text generated by n-gram models trained on a novel (Mary Shelley's
\emph{Frankenstein}), for increasing $n$:
\begin{center}
\small
\begin{tabularx}{\linewidth}{@{}c X@{}}
\toprule
$n$ & generated text \\
\midrule
1 & ``care I of . its destined , in , from . . for the Felix an not which measure excited'' \\
2 & ``I compassionated him , how heavily ; but I had arrived , but they averred , my unhallowed damps and'' \\
3 & ``I continued their single offspring . At length I gathered from a man who , born in freedom , spurned'' \\
4 & ``I continued walking in this manner , during which I enjoyed the feeling of happiness . Still thou canst listen'' \\
5 & ``I continued walking in this manner for some time , and I feared the effects of the dæmon's disappointment .'' \\
\bottomrule
\end{tabularx}
\end{center}
\begin{itemize}[nosep]
  \item Increasing $n$, the model is \textbf{more coherent}: unigrams give a word salad,
        bigrams give locally plausible pairs, 4- and 5-grams read like real sentences.
  \item Increasing $n$ \textbf{increases the sparsity problem}.
  \item With a high $n$ the model \textbf{repeats the training set}: most 5-word
        contexts were seen only once, so they have only one possible continuation, and
        the model copies whole passages of the novel instead of creating new text.
\end{itemize}

\begin{nbsummary}[n-gram language models]
\begin{itemize}[nosep]
  \item Markov assumption: the next word depends only on the previous $n-1$ words.
  \item Estimate: $p(x \mid \text{context}) \approx
        \text{count}(\text{context}, x) / \text{count}(\text{context})$.
  \item Sparsity 1 (unseen n-gram, probability 0) $\to$ smoothing; sparsity 2 (unseen
        context, division by 0) $\to$ backoff. Both get worse as $n$ grows; $n \le 5$.
  \item Storage grows with $n$ and with the corpus.
  \item Small $n$: incoherent text; large $n$: coherent but copied from the corpus.
  \item The context is cut off: anything older than $n-1$ words is forgotten.
\end{itemize}
\end{nbsummary}

\begin{nbexercise}{1}{bigram probabilities}
\notinslides
Using the toy corpus of this section (\texttt{<s>} I like cats \texttt{</s>},
\texttt{<s>} I like dogs \texttt{</s>}, \texttt{<s>} you like cats \texttt{</s>}) and a
bigram model without smoothing, compute $p(\textit{you like cats})$. Which sentence of
the form ``\emph{subject} like \emph{animal}'' has the highest probability?
\end{nbexercise}
\begin{nbsolution}
Counts: $p(\textit{you} \mid \texttt{<s>}) = \frac13$,
$p(\textit{like} \mid \textit{you}) = \frac11 = 1$,
$p(\textit{cats} \mid \textit{like}) = \frac23$,
$p(\texttt{</s>} \mid \textit{cats}) = \frac22 = 1$. So
\[
  p(\textit{you like cats}) = \tfrac13 \cdot 1 \cdot \tfrac23 \cdot 1 = \tfrac29 .
\]
The four candidates are
$p(\textit{I like cats}) = \frac23 \cdot \frac23 = \frac49$,
$p(\textit{I like dogs}) = \frac23 \cdot \frac13 = \frac29$,
$p(\textit{you like cats}) = \frac13 \cdot \frac23 = \frac29$,
$p(\textit{you like dogs}) = \frac13 \cdot \frac13 = \frac19$ (the factors equal to 1
are omitted). Check: they add up to $\frac{4+2+2+1}{9} = 1$, since these are the only
sentences the model can produce. \emph{I like cats} wins: it combines the most
frequent start (\emph{I}) and the most frequent animal (\emph{cats}).
\end{nbsolution}

\subsection{Neural Language Models}
\label{sec:neural-language-models}
\slidesrc{Language models \& RNNs: 35--50}

n-gram models fail because they treat every word sequence as a separate, unrelated
item: seeing ``she went to school'' a thousand times teaches them nothing about ``he
walked to college''. Neural language models fix this by representing words as
\textbf{vectors}, so that similar words get similar representations and what is learnt
about one word transfers to its neighbours.

\subsubsection{Aside: word embeddings via word2vec}
\slidesrc{Language models \& RNNs: 36}

\paragraph{One-hot vectors.}
\notinslides
A computer needs numbers, not words. The simplest way to turn a word into numbers is a
\textbf{one-hot vector}: number the words of the vocabulary from 1 to $|V|$; word number
$i$ becomes a vector of length $|V|$ with a 1 in position $i$ and 0 everywhere else.
With the vocabulary $(\textit{cat}, \textit{dog}, \textit{the}, \textit{sleeps})$:
\[
  \textit{cat} = \col{1\\0\\0\\0}, \quad \textit{dog} = \col{0\\1\\0\\0}, \quad
  \textit{the} = \col{0\\0\\1\\0}, \quad \textit{sleeps} = \col{0\\0\\0\\1}.
\]
One-hot vectors identify words but say nothing about their meaning: \emph{cat} is
exactly as different from \emph{dog} as from \emph{sleeps} (any two of them have
dot product 0).

\paragraph{Word embeddings.}
The goal of a \textbf{word embedding} is to represent each word with a short, dense
vector (a few hundred real numbers) such that:
\begin{itemize}[nosep]
  \item the vector of a word is \textbf{similar} to the vectors of words that occur in
        \emph{similar} contexts;
  \item it is \textbf{dissimilar} from the vectors of words that are unlikely to appear
        in the same contexts.
\end{itemize}
This is the distributional hypothesis (Section~\ref{sec:distributional-hypothesis})
turned into vectors (Section~\ref{sec:geometric-interpretation}). \textbf{Word2vec} is
a popular method that learns such vectors by training a small neural network to
predict words from their neighbours. Similarity between two vectors is measured by
their \textbf{dot product}: multiply the entries position by position and add.

\begin{nbexample}{similar and dissimilar words}
\slidesrc{Language models \& RNNs: 36}
Eight-dimensional embeddings from the slide:
\[
  \textit{baking} = \col{0.286\\0.792\\-0.177\\-0.107\\0.109\\-0.542\\0.349\\0.271}\!,\quad
  \textit{flour}  = \col{0.413\\0.582\\-0.007\\0.247\\0.216\\-0.718\\0.147\\0.051}\!,\quad
  \textit{dollar} = \col{-0.154\\0.024\\0.512\\0.782\\0.359\\0.141\\-0.581\\-0.027}\!.
\]
\begin{align*}
  \textit{baking}\T\textit{flour}
  &= 0.286 \cdot 0.413 + 0.792 \cdot 0.582 + (-0.177)(-0.007) + \cdots\\
  &\qquad + 0.271 \cdot 0.051\\
  &= 0.118 + 0.461 + 0.001 - 0.026 + 0.024 + 0.389 + 0.051 + 0.014\\
  &\approx 1.03
\end{align*}
The large positive entries are in the same places (positions 2 and 6 contribute most),
so the dot product is large: \emph{baking} and \emph{flour} appear in similar contexts
(recipes). For \emph{baking} and \emph{dollar} the signs disagree in most positions and
the dot product is negative: they are used in unrelated contexts.
\end{nbexample}

\begin{nbwarning}[the number on the slide]
The slide gives $\textit{baking}\T\textit{dollar} = -0.35$. With the vectors as printed
the dot product is $-0.45$; $-0.34$ is their \emph{cosine similarity} (the dot product
divided by the lengths of the two vectors, Section~\ref{sec:geometric-interpretation}).
Either way the conclusion is the same: negative, hence dissimilar.
\end{nbwarning}

\subsubsection{A fixed-window neural language model}
\slidesrc{Language models \& RNNs: 37--41}

Recall the language-modelling task: given a sequence of words
$\vect x^{(1)}, \ldots, \vect x^{(t)}$, output the probability distribution of the next
word, $p(\vect x^{(t+1)} \mid \vect x^{(1)}, \ldots, \vect x^{(t)})$. The first neural
approach keeps the n-gram idea of a \textbf{fixed window} of previous words, but
replaces counting with a \textbf{neural network} that \emph{learns} to predict.

For ``Anna was tired so she went to her \underline{\quad}'', a window of 4 words
takes \emph{she went to her} and predicts the missing word. The network is built in
four layers, from bottom to top:

\begin{center}
\begin{tikzpicture}[>=Stealth,
    neu/.style={circle, draw=#1, fill=#1!35, minimum size=5pt, inner sep=0pt},
    lay/.style={draw=#1, rounded corners, thick, inner sep=2pt},
    lab/.style={font=\small, align=left, anchor=west, text width=4.3cm}]
  % one-hot words
  \foreach \w [count=\i] in {she, went, to, her} {
    \node[font=\small] (w\i) at (1.6*\i,0) {\w};
    \node[font=\scriptsize, nbgray] at (1.6*\i,-0.4) {$\vect x^{(\i)}$};
    % embedding block
    \foreach \k in {1,...,4} \node[neu=magenta] (e\i\k) at (1.6*\i+0.24*\k-0.6,1.2) {};
    \draw[->, nbgray] (w\i) -- node[right, font=\scriptsize] {$E$} (1.6*\i,1.05);
  }
  \node[lay=magenta, fit=(e11)(e44)] (emb) {};
  % hidden
  \foreach \k in {1,...,10} \node[neu=cyan] (h\k) at (2.9+0.32*\k,2.5) {};
  \node[lay=cyan, fit=(h1)(h10)] (hid) {};
  \draw[->, thick] (emb) -- node[right, font=\scriptsize] {$W_e$} (hid);
  % output histogram
  \node[lay=nbred, minimum width=5.6cm, minimum height=0.45cm, font=\scriptsize]
       (out) at (4.65,3.8) {distribution over the whole vocabulary};
  \draw[->, thick] (hid) -- node[right, font=\scriptsize] {$U$} (out);
  \foreach \x/\h in {0.2/0.15,0.5/0.4,0.8/0.2,1.3/0.6,1.6/0.25,2.1/0.3,2.6/0.1,
                     3.0/0.8,3.3/0.35,3.9/0.2,4.4/0.12,4.9/0.3}
    \fill[nbred!70] (1.85+\x,4.05) rectangle ++(0.15,\h);
  % labels
  \node[lab, nbred]   at (8.0,3.85) {\textbf{output}\\$\hat{\vect y} = \softmax(U\vect h) \in \mathbb R^{|V|}$};
  \node[lab, cyan!60!black] at (8.0,2.5) {\textbf{hidden layer}\\$\vect h = \sigma(W_e \vect e)$};
  \node[lab, magenta] at (8.0,1.2) {\textbf{word embeddings}\\$\vect e = [\vect e^{(1)}, \vect e^{(2)}, \vect e^{(3)}, \vect e^{(4)}]$};
  \node[lab]          at (8.0,-0.15) {\textbf{one-hot word vectors}\\$\vect x^{(1)}, \vect x^{(2)}, \vect x^{(3)}, \vect x^{(4)}$};
\end{tikzpicture}
\captionof{figure}{The fixed-window neural language model. Information flows upwards:
four words $\to$ their embeddings, glued together $\to$ a hidden layer $\to$ a
probability for every word of the vocabulary (tall bars = likely next words).}
\end{center}

\begin{enumerate}
  \item \textbf{One-hot word vectors} $\vect x^{(1)}, \ldots, \vect x^{(4)} \in
        \mathbb R^{|V|}$: the four words of the window.
  \item \textbf{Word embeddings.} Each word is turned into its embedding
        $\vect e^{(i)} = E \vect x^{(i)}$. $E$ is a matrix with one column per word of
        the vocabulary ($E \in \mathbb R^{d \times |V|}$, $d$ = embedding size).
        Multiplying a matrix by a one-hot vector simply \emph{picks one column}: the
        column of that word. So $E$ is just a table of word vectors, stored as a
        matrix. The four embeddings are then \textbf{concatenated} (written one after
        the other) into a single vector $\vect e$ of length $4d$.
  \item \textbf{Hidden layer} $\vect h = \sigma(W_e \vect e)$. The matrix $W_e$ mixes all
        the numbers of the four embeddings; $\sigma$ is a non-linear function applied
        to each entry (see below). $\vect h$ is a summary of the window useful for
        prediction.
  \item \textbf{Output} $\hat{\vect y} = \softmax(U \vect h) \in \mathbb R^{|V|}$: one
        number per word of the vocabulary, turned into a probability distribution by
        the softmax (see below). $\hat{\vect y}_w$ is the predicted probability that
        the next word is $w$. The hat in $\hat{\vect y}$ means ``predicted''.
\end{enumerate}

\notinslides
Two building blocks deserve a closer look.
\begin{itemize}
  \item \textbf{The non-linearity $\sigma$.} Typically the \emph{sigmoid}
        $\sigma(z) = \frac{1}{1 + e^{-z}}$, which squashes any number into $(0, 1)$, or
        $\tanh$, which squashes into $(-1, 1)$. Without it, the whole network would be
        a product of matrices, i.e.\ just one big linear map, and could only learn
        very simple relations. (Real networks also add a \emph{bias} vector inside
        $\sigma$, $\sigma(W_e \vect e + \vect b)$; the slides omit it for simplicity.)
  \item \textbf{The weights $E$, $W_e$, $U$} are the \textbf{parameters} of the model,
        written together as $\theta$. They start as random numbers; \emph{training}
        (below) adjusts them so that the predictions become good.
\end{itemize}

\subsubsection{How to train a neural language model}
\slidesrc{Language models \& RNNs: 42--49}

\begin{nbdef}{The four ingredients of training a neural network}
\begin{enumerate}[nosep]
  \item \textbf{A corpus of data}: here a collection of documents, i.e.\ a long sequence
        of words $\vect x^{(1)}, \ldots, \vect x^{(T)}$.
  \item \textbf{A model}: here the fixed-window neural LM.
  \item \textbf{A loss function $J$}: a number measuring how wrong the model's
        predictions are; it compares the output distribution with the true one.
  \item \textbf{An optimisation algorithm} that changes the parameters to make the
        loss smaller: \textbf{stochastic gradient descent}.
\end{enumerate}
\end{nbdef}

\paragraph{The softmax.}
The vector $U\vect h \in \mathbb R^{|V|}$ contains one \emph{score} per word, but scores
can be negative and do not add up to 1: it is an \textbf{un-normalised} probability
distribution. The \textbf{softmax} normalises it:
\[
  \softmax(\vect z)_i = \frac{e^{z_i}}{\sum_{j=1}^{K} e^{z_j}}, \qquad i = 1, \ldots, K .
\]
The exponential makes every entry positive (and enlarges the gaps between scores),
and dividing by the sum makes the entries add up to 1.

\begin{nbexample}{softmax}
\slidesrc{Language models \& RNNs: 47}
Scores $U\vect h = (1.3,\ 5.1,\ 2.2,\ 0.7,\ 1.1)$ for a vocabulary of $K = 5$ words.
\begin{align*}
  (e^{1.3}, e^{5.1}, e^{2.2}, e^{0.7}, e^{1.1})
  &= (3.67,\ 164.02,\ 9.03,\ 2.01,\ 3.00) \reason{exponentiate}\\
  \text{sum} &= 3.67 + 164.02 + 9.03 + 2.01 + 3.00 = 181.73 \\
  \hat{\vect y} &= (0.02,\ 0.90,\ 0.05,\ 0.01,\ 0.02) \reason{divide by the sum}
\end{align*}
Check: $0.02 + 0.90 + 0.05 + 0.01 + 0.02 = 1$. The highest score (5.1) gets most of the
probability, but no word gets exactly 0.
\end{nbexample}

\paragraph{The loss: cross-entropy.}
We compare
\begin{itemize}[nosep]
  \item the \textbf{computed} probability $\hat y_i = q(V_i \mid \textit{she went to
        her})$, what the model says about word $V_i$ of the vocabulary;
  \item the \textbf{true} probability $y^*_i = p(V_i \mid \text{full context})$, what
        the word really is.
\end{itemize}
The standard way to compare two distributions is the \textbf{cross-entropy}:
\[
  \mathrm{CE}(\vect y^*, \hat{\vect y}) = - \sum_{i=1}^{|V|} y^*_i \log \hat y_i .
\]
In training data we do not know the ``true distribution'': we only know which word
\emph{actually came next} in the corpus. So $\vect y^*$ is a \textbf{one-hot vector}:
1 for the true next word, 0 for every other word. All the terms of the sum vanish
except one. If the true word in the data is the fifth word of the sentence,
$\vect x^{(5)}$, then
\[
  J^{(t)} = \mathrm{CE}(\vect y^*, \hat{\vect y}) = -\log \hat y_{\vect x^{(5)}}
  \qquad \text{(the negative log-probability of the true word).}
\]
(The slide writes $\hat y_5$, meaning ``the entry of $\hat{\vect y}$ belonging to the
true word $\vect x^{(5)}$''.)

\begin{nbintuition}[why $-\log$ of the probability?]
\notinslides
$-\log p$ is 0 when $p = 1$ (the model was certain and right: no penalty) and grows
without limit as $p \to 0$ (the model said the true word was nearly impossible: huge
penalty). With the softmax output above (natural logarithm):
\begin{itemize}[nosep]
  \item if the true word is the 2nd one ($\hat y = 0.90$): loss
        $= -\ln 0.90 \approx 0.10$, small;
  \item if the true word is the 4th one ($\hat y = 0.011$): loss
        $= -\ln 0.011 \approx 4.50$, large.
\end{itemize}
Minimising the loss means: \emph{give high probability to the words that really come
next}.
\end{nbintuition}

\paragraph{Gradient descent.}
\notinslides
The loss $J(\theta)$ depends on all the parameters $\theta$ (the entries of $E$, $W_e$,
$U$). The \textbf{gradient} $\nabla_\theta J$ is the vector of the derivatives of $J$
with respect to each parameter: it points in the direction in which the loss
\emph{increases} fastest. \textbf{Gradient descent} takes a small step in the opposite
direction, again and again:
\[
  \theta \leftarrow \theta - \alpha\, \nabla_\theta J(\theta),
\]
where $\alpha > 0$ (the \textbf{learning rate}) controls the step size. It is like
walking down a hill in the fog: you cannot see the valley, but you can feel the slope
under your feet and step downhill. The gradients are computed automatically by
\textbf{backpropagation}, i.e.\ the chain rule of calculus applied layer by layer from
the output back to the input. \emph{Stochastic} gradient descent computes the gradient
on a small random part of the data at each step instead of on the whole corpus
(Section~\ref{sec:recurrent-neural-networks}).

\subsubsection{Pros and cons of the fixed-window model}
\slidesrc{Language models \& RNNs: 50}

This model was introduced by Y.~Bengio et al., \emph{A Neural Probabilistic Language
Model}, 2003.

\textbf{Improvements over n-gram LMs:}
\begin{itemize}[nosep]
  \item \textbf{no sparsity problem}: the network computes a probability for every word
        from the embeddings, so unseen sequences get sensible, non-zero probabilities;
        if \emph{school} and \emph{college} have similar embeddings, what is learnt
        for one transfers to the other;
  \item \textbf{no need to store n-grams}: only the weights are stored.
\end{itemize}

\textbf{Remaining problems:}
\begin{itemize}[nosep]
  \item the \textbf{fixed window is too restrictive}: like an n-gram model it cannot see
        anything older than its window (``Anna was \emph{tired}'' is still invisible);
  \item \textbf{enlarging the window increases the size of $W_e$}: $W_e$ has one column
        per entry of $\vect e$, and $\vect e$ grows with every extra word; a window big
        enough for any text is impossible;
  \item \textbf{no symmetry}: $\vect x^{(1)}$ and $\vect x^{(2)}$ are multiplied by
        \emph{different} columns of $W_e$. What the model learns about the word
        \emph{school} in position 4 does not transfer to \emph{school} in position 2:
        the same knowledge has to be learnt separately for every position.
\end{itemize}

\begin{nbwhy}
The question left open by this slide, \emph{how can we process inputs of arbitrary
length?}, is answered by recurrent neural networks (next section), and later by
transformers. Both keep the two good ideas introduced here: words as learnt
embeddings, and a softmax over the vocabulary trained with cross-entropy. Every modern
LLM, GPT included, still ends with exactly this softmax layer.
\end{nbwhy}

\begin{nbsummary}[neural language models]
\begin{itemize}[nosep]
  \item Embedding: $\vect e = E\vect x$ (picks the column of word $\vect x$); similar
        words have a large dot product.
  \item Fixed window: $\hat{\vect y} = \softmax(U\, \sigma(W_e [\vect e^{(1)}, \ldots,
        \vect e^{(n)}]))$.
  \item $\softmax(\vect z)_i = e^{z_i} / \sum_j e^{z_j}$: positive, sums to 1.
  \item Loss: cross-entropy with a one-hot target $= -\log \hat y_{\text{true word}}$.
  \item Training: corpus + model + loss + (stochastic) gradient descent
        $\theta \leftarrow \theta - \alpha \nabla_\theta J$.
  \item Solves sparsity and storage; still a fixed window, growing $W_e$, no weight
        sharing across positions.
\end{itemize}
\end{nbsummary}

\subsection{Recurrent Neural Networks}
\label{sec:recurrent-neural-networks}
\slidesrc{Language models \& RNNs: 51--84}

\subsubsection{The idea: reuse the same weights at every step}
\slidesrc{Language models \& RNNs: 52--56}

A \textbf{recurrent neural network} (RNN) reads a text \emph{one word at a time}, the
way we do, and keeps a \textbf{hidden state} $\vect h^{(t)}$: a vector that summarises
everything read so far. At each step it combines the previous summary with the new
word to produce the new summary. Crucially, \textbf{the same weights are applied
repeatedly} at every step, so the network has the same size whether the text has 4
words or 4000: inputs can be of \textbf{any length}.

\begin{center}
\begin{tikzpicture}[>=Stealth,
    neu/.style={circle, draw=cyan!70!black, fill=cyan!35, minimum size=5pt, inner sep=0pt},
    hid/.style={draw=cyan!70!black, rounded corners, thick, minimum width=0.5cm,
                minimum height=1.45cm}]
  \foreach \w [count=\i] in {the, quick, brown, fox} {
    \pgfmathsetmacro{\x}{2.3*\i}
    \node[hid] (h\i) at (\x,0) {};
    \foreach \k in {1,...,4} \node[neu] at (\x,0.72-0.29*\k) {};
    \node[font=\scriptsize, anchor=south west, inner sep=1pt] at (h\i.north east) {$\vect h^{(\i)}$};
    \node[font=\small] (x\i) at (\x,-1.7) {\w};
    \node[font=\scriptsize, nbgray] at (\x,-2.05) {$\vect x^{(\i)}$};
    \node[font=\small] (y\i) at (\x,1.7) {$\hat{\vect y}^{(\i)}$};
    \draw[->, thick, nbgray] (x\i) -- (h\i);
    \draw[->, thick, nbgray] (h\i) -- (y\i);
  }
  \node[hid, draw=nbgray, minimum height=1.45cm] (h0) at (0,0) {};
  \node[font=\scriptsize, anchor=south west, inner sep=1pt] at (h0.north east) {$\vect h^{(0)}$};
  \foreach \i/\j in {0/1,1/2,2/3,3/4} {
    \draw[->, thick, nbgray] (h\i) -- node[draw=nbred, fill=white, font=\scriptsize,
        inner sep=2pt] (w\j) {$W_h$} (h\j);
  }
  \draw[->, thick, nbgray] (h4) -- ++(1.0,0) node[right] {\ldots};
  \node[font=\footnotesize, align=right, anchor=east] at (-0.6,1.7) {outputs};
  \node[font=\footnotesize, align=right, anchor=east] at (-0.6,0) {hidden\\states};
  \node[font=\footnotesize, align=right, anchor=east] at (-0.6,-1.85) {inputs, of\\any length};
  \node[draw=nbred, rounded corners, fill=nbred!5, font=\footnotesize, align=center] at (5.75,-2.9)
       {all the red boxes are the \emph{same} weights $W_h$,\\applied repeatedly to the whole sequence};
\end{tikzpicture}
\captionof{figure}{An RNN ``unrolled'' over four time steps. There is only one matrix
$W_h$ (and one $W_e$, one $U$): the copies are the same weights used again at every
step.}
\end{center}

\begin{nbdef}{RNN language model}
At each time step $t = 1, 2, \ldots$
\begin{align*}
  \vect e^{(t)} &= E \vect x^{(t)}
      && \text{word embedding of the one-hot word } \vect x^{(t)} \in \mathbb R^{|V|}\\
  \vect h^{(t)} &= \sigma\bigl(W_h \vect h^{(t-1)} + W_e \vect e^{(t)}\bigr)
      && \text{new hidden state} \\
  \hat{\vect y}^{(t)} &= \softmax\bigl(U \vect h^{(t)}\bigr) \in \mathbb R^{|V|}
      && \text{distribution of the next word } \vect x^{(t+1)}
\end{align*}
$\vect h^{(0)}$ is an initial hidden state (often a vector of zeros). The parameters are
$E$, $W_h$, $W_e$, $U$, shared by all time steps.
\end{nbdef}

Read the hidden-state equation slowly: $W_h \vect h^{(t-1)}$ is ``what I remember'',
$W_e \vect e^{(t)}$ is ``what I just read''; they are added and squashed by $\sigma$ into
the new memory $\vect h^{(t)}$. Since $\vect h^{(t-1)}$ itself depended on
$\vect h^{(t-2)}$ and on $\vect x^{(t-1)}$, and so on, $\vect h^{(t)}$ depends, in
principle, on \emph{all} the words $\vect x^{(1)}, \ldots, \vect x^{(t)}$. There is no
window any more.

\begin{nbintuition}[an RNN is a reader taking notes]
\notinslides
Imagine reading a novel with a notebook of fixed size. After each word you update your
notes (same rules every time), and from the notes alone you guess the next word. The
notebook never grows, whatever the length of the book, so you must decide what to keep
and what to overwrite. This is both the strength of RNNs (any length, constant size)
and, as we will see, their weakness (old information gets overwritten).
\end{nbintuition}

\needspace{12\baselineskip}
\subsubsection{Pros and cons}
\slidesrc{Language models \& RNNs: 57}

\begin{center}
\small
\begin{tabularx}{\linewidth}{@{}X X@{}}
\toprule
\textbf{Pros} & \textbf{Cons} \\
\midrule
not restricted to fixed-length inputs &
  computation is slow: step $t$ needs $\vect h^{(t-1)}$, so the steps cannot be computed
  in parallel \\
can use information from many time steps ago (in theory) &
  in practice long-range information is hard to capture
  (Section~\ref{sec:vanishing-exploding-gradients}) \\
input words are processed symmetrically: the same weights at every position
(\textbf{weight sharing}), so what is learnt about a word in one position applies
everywhere & \\
\bottomrule
\end{tabularx}
\end{center}
Compare with the fixed-window model (Section~\ref{sec:neural-language-models}): the
first and third pros solve exactly its ``fixed window'' and ``no symmetry'' problems,
and the model size no longer depends on the length of the context.

\subsubsection{Training an RNN language model}
\slidesrc{Language models \& RNNs: 58--67}

\begin{enumerate}
  \item Gather a large corpus of text $\vect x^{(1)}, \ldots, \vect x^{(T)}$.
  \item Feed it to the RNN one word at a time, computing $\hat{\vect y}^{(t)}$ for every
        $t = 1, \ldots, T$.
  \item At each step compute the loss: the cross-entropy between the predicted
        distribution $\hat{\vect y}^{(t)}$ and the true next word
        $\vect y^{(t)} = \vect x^{(t+1)}$ (one-hot):
        \[
          J^{(t)} = \mathrm{CE}(\vect y^{(t)}, \hat{\vect y}^{(t)})
          = -\sum_{w \in V} \underbrace{y^{(t)}_w}_{\text{one-hot}} \log \hat y^{(t)}_w
          = -\log \hat y^{(t)}_{\vect x^{(t+1)}} ,
        \]
        the negative log-probability of the correct next word.
  \item Average the loss over the whole dataset (the last word has no next word to
        predict, hence $T-1$ terms):
        \[
          J = \frac{1}{T-1} \sum_{t=1}^{T-1} J^{(t)}
            = \frac{1}{T-1} \sum_{t=1}^{T-1} -\log \hat y^{(t)}_{\vect x^{(t+1)}} .
        \]
\end{enumerate}

\begin{center}
\begin{tikzpicture}[>=Stealth,
    neu/.style={circle, draw=#1!70!black, fill=#1!35, minimum size=4pt, inner sep=0pt},
    box/.style={draw=#1!70!black, rounded corners, minimum width=0.42cm, minimum height=1.15cm}]
  \foreach \w/\nx [count=\i] in {the/quick, quick/brown, brown/fox, fox/jumped} {
    \pgfmathsetmacro{\x}{2.2*\i}
    \node[box=magenta] (e\i) at (\x,0) {};
    \foreach \k in {1,...,4} \node[neu=magenta] at (\x,0.57-0.23*\k) {};
    \node[box=cyan] (h\i) at (\x,1.9) {};
    \foreach \k in {1,...,4} \node[neu=cyan] at (\x,2.47-0.23*\k) {};
    \node[font=\small] (x\i) at (\x,-1.3) {\w};
    \node[font=\scriptsize] (y\i) at (\x,3.25) {$\hat{\vect y}^{(\i)}$};
    \node[font=\scriptsize, nbred] (J\i) at (\x,4.05) {$J^{(\i)}$};
    \draw[->, nbgray] (x\i) -- node[right, font=\tiny] {$E$} (e\i);
    \draw[->, nbgray] (e\i) -- node[right, font=\tiny] {$W_e$} (h\i);
    \draw[->, nbgray] (h\i) -- node[right, font=\tiny] {$U$} (y\i);
    \draw[->, nbred] (y\i) -- (J\i);
    \node[font=\tiny, nbred, align=center] at (\x,4.55) {$-\log \hat y^{(\i)}_{\textit{\nx}}$};
  }
  \node[box=cyan, draw=nbgray] (h0) at (0,1.9) {};
  \node[font=\scriptsize] at (0,2.75) {$\vect h^{(0)}$};
  \foreach \i/\j in {0/1,1/2,2/3,3/4} \draw[->, nbgray] (h\i) -- node[above, font=\tiny] {$W_h$} (h\j);
  \draw[->, nbgray] (h4) -- ++(0.8,0) node[right] {\ldots};
  \node[font=\small] at (11.0,-1.3) {jumped};
  \node[font=\scriptsize, nbred, anchor=east] at (-0.3,4.05) {loss at each step};
  \node[font=\scriptsize, nbred, anchor=east] at (-0.3,3.25) {predicted distributions};
  \node[font=\scriptsize, nbred, anchor=east] at (-0.3,-1.3) {corpus};
  \node[font=\small, nbred] at (5.5,5.2)
       {$J = \frac14\bigl(J^{(1)} + J^{(2)} + J^{(3)} + J^{(4)}\bigr)$};
\end{tikzpicture}
\captionof{figure}{Training on ``the quick brown fox jumped''. At step 1 the network has
read \emph{the} and is penalised by $-\log$ of the probability it gave to
\emph{quick}; at step 2 by the probability it gave to \emph{brown}; and so on. The total
loss is the average.}
\end{center}

Note that the input at every step is the \emph{true} word from the corpus, not the
word the model predicted: during training the model is always shown the correct text
and asked to predict one step ahead.

\paragraph{Training in practice: stochastic gradient descent.}
A typical corpus has millions of words. Computing the loss
$J(\theta) = \frac{1}{T-1}\sum_{t} J^{(t)}(\theta)$ and its gradient over all of it at
once requires too much memory and time. So:
\begin{itemize}[nosep]
  \item divide the corpus into smaller \textbf{chunks} (sentences, documents);
  \item \textbf{Stochastic Gradient Descent} (SGD) approximates the full gradient with
        the average gradient over a small \textbf{batch} of $S$ chunks:
        \[
          \nabla_\theta J(\theta) \approx \frac1S \sum_{j=1}^{S} \nabla_\theta J^{(j)}(\theta) ;
        \]
  \item loss and gradients are computed for one batch of sentences at a time, then the
        weights are updated ($\theta \leftarrow \theta - \alpha \nabla_\theta J$) and the
        next batch is taken.
\end{itemize}
The approximation is noisy (a different batch gives a slightly different gradient),
but on average it points in the right direction, and each step is cheap. That noise is
the ``stochastic'' in the name.

\subsubsection{Backpropagation through time}
\slidesrc{Language models \& RNNs: 68--76}

To run gradient descent we need the gradient of the loss with respect to every
parameter. The tricky one is $W_h$, because it is used at \emph{every} step. To see the
structure clearly, the slides use an even simpler model with no input and no
non-linearity:
\[
  \vect h^{(t)} = W_h \vect h^{(t-1)}, \qquad
  J^{(t)}(\theta) = \mathrm{CE}(\vect h^{(t)}, \vect y^{*(t)}) .
\]

\begin{nbfact}{Backpropagation through time (BPTT)}
The gradient with respect to a repeated weight is the \textbf{sum of the gradients for
each time it appears}:
\[
  \frac{\partial J^{(t)}}{\partial W_h} =
  \sum_{i=1}^{t} \left. \frac{\partial J^{(t)}}{\partial W_h} \right|_{(i)},
\]
where $\bigl.\frac{\partial J^{(t)}}{\partial W_h}\bigr|_{(i)}$ is the gradient obtained by
treating only the copy of $W_h$ used at step $i$ as variable.
\end{nbfact}

\begin{nbintuition}[why a sum?]
\notinslides
Think of $W_h$ as a knob that turns $t$ dials at once (one per step). Turning it a
little changes the output through each dial; for small changes the effects simply add
up. So the total sensitivity is the sum of the sensitivities through each use.
\end{nbintuition}

\paragraph{The derivation.}
Start with the \textbf{chain rule}: the loss depends on $W_h$ only through
$\vect h^{(t)}$, so
\[
  \frac{\partial J^{(t)}}{\partial W_h} =
  \frac{\partial J^{(t)}}{\partial \vect h^{(t)}}\,
  \frac{\partial \vect h^{(t)}}{\partial W_h}.
\]
The first factor is easy (it only involves the output). The second is the interesting
one: $\vect h^{(t)} = W_h \vect h^{(t-1)}$ depends on $W_h$ in two ways, \emph{directly}
(it is multiplied by $W_h$) and \emph{indirectly} (through $\vect h^{(t-1)}$, which was
itself computed with $W_h$). Write it as $\vect h^{(t)}(W_h, \vect h^{(t-1)})$.

To keep the algebra readable we first do it with \textbf{numbers instead of matrices}
(a hidden state of size 1, so $h^{(t)} = w\, h^{(t-1)}$). The derivative of a product is
``derivative of the first times the second plus the first times the derivative of the
second'' (product rule), so:
\begin{align*}
  \frac{\dd h^{(t)}}{\dd w}
  &= \underbrace{h^{(t-1)}}_{\text{direct}}
     + \underbrace{w\, \frac{\dd h^{(t-1)}}{\dd w}}_{\text{indirect}}
     \reason{product rule on $w\,h^{(t-1)}$} \\
  &= h^{(t-1)} + w \Bigl( h^{(t-2)} + w\, \frac{\dd h^{(t-2)}}{\dd w} \Bigr)
     \reason{same rule for $h^{(t-1)}$} \\
  &= h^{(t-1)} + w\, h^{(t-2)} + w^2\, \frac{\dd h^{(t-2)}}{\dd w}
     \reason{expand} \\
  &= \cdots = h^{(t-1)} + w\, h^{(t-2)} + w^2 h^{(t-3)} + \cdots + w^{t-1} h^{(0)}
     \reason{$\dd h^{(0)}/\dd w = 0$} \\
  &= \sum_{i=0}^{t-1} w^{\,t-1-i}\, h^{(i)} .
\end{align*}
Each term is one ``use'' of $w$: the term with $h^{(i)}$ comes from the copy of $w$
that multiplies $h^{(i)}$ at step $i+1$, and the factor $w^{t-1-i}$ carries its effect
through the remaining $t-1-i$ steps up to $h^{(t)}$. This is the BPTT sum.

\begin{nbexample}{checking the formula}
\notinslides
Take $w = 0.7$, $h^{(0)} = 2$, $t = 5$. Since $h^{(t)} = w^t h^{(0)}$, directly
$\frac{\dd h^{(5)}}{\dd w} = 5 w^4 h^{(0)} = 5 \cdot 0.2401 \cdot 2 = 2.401$. With the
formula, $h^{(i)} = 2 \cdot 0.7^i$ and every term equals $w^{4-i} \cdot 2 w^i = 2w^4$:
\[
  \sum_{i=0}^{4} w^{4-i} h^{(i)} = 5 \cdot 2 \cdot 0.7^4 = 2.401 . \quad\checkmark
\]
\end{nbexample}

With matrices the computation is the same, and the slides write the result as
\[
  \frac{\partial \vect h^{(t)}}{\partial W_h}
  = \sum_{i=0}^{t-1} \Bigl( \prod_{j=i+1}^{t-1} W_h \Bigr) \vect h^{(i)}
  = \sum_{i=0}^{t-1} W_h^{\,t-1-i}\, \vect h^{(i)} .
\]

\begin{nbwarning}[what the matrix formula means]
The derivative of a vector with respect to a matrix is a three-index object, so the
slide's formula (and its intermediate step ``$\vect h^{(t-1)} I$'') is a shorthand. The
precise statement is: if $W_h$ is changed by a small matrix $\Delta W$, then
$\vect h^{(t)}$ changes by
\[
  \Delta \vect h^{(t)} \approx \sum_{i=0}^{t-1} W_h^{\,t-1-i}\, \Delta W\, \vect h^{(i)} .
\]
The structure is exactly the scalar one: one term per use of $W_h$, each carried to step
$t$ by a power of $W_h$.
\end{nbwarning}

\begin{itemize}[nosep]
  \item BPTT is \textbf{equivalent to ordinary backpropagation through the whole
        unrolled network}: once the RNN is drawn as a long chain of copies, it is just a
        (very deep) feed-forward network whose copies share the same weights.
  \item In practice the backward pass is \textbf{truncated} at about 20 steps, for
        efficiency: terms older than that are ignored (\textbf{truncated BPTT}).
  \item For more complex RNNs (with inputs and non-linearities) the formula is very
        similar.
\end{itemize}

\subsubsection{Generating text with an RNN}
\slidesrc{Language models \& RNNs: 77}

To generate text at test time, \textbf{repeatedly sample} from the network,
\textbf{feeding each output back as the next input}:
\begin{enumerate}[nosep]
  \item feed the start token \texttt{<s>}; the network outputs $\hat{\vect y}^{(1)}$;
  \item sample a word from $\hat{\vect y}^{(1)}$, say \emph{the};
  \item feed \emph{the} as the next input; the network updates its hidden state and
        outputs $\hat{\vect y}^{(2)}$; sample \emph{quick};
  \item feed \emph{quick}, sample \emph{brown}, and so on, until \texttt{</s>} is
        sampled or a length limit is reached.
\end{enumerate}
Unlike an n-gram model, the hidden state carries information from \emph{all} the words
generated so far, not only the last $n-1$. This is also, in essence, how chatbots
write their answers, one token at a time.

\subsubsection{Evaluating language models: perplexity}
\slidesrc{Language models \& RNNs: 78--84}

How do we measure whether a language model is good? We take a \textbf{test sentence}
$X_{\text{test}} = \vect x^{(1)}, \ldots, \vect x^{(T)}$ that was \textbf{not seen
during training} (otherwise a model that memorised the corpus would look perfect) and
check how much probability the model gives it.

\begin{nbdef}{Perplexity}
The \textbf{perplexity} of a sentence is its inverse probability, normalised by its
length $T$:
\[
  \text{perplexity}(X_{\text{test}}) = p(\vect x^{(1)}, \ldots, \vect x^{(T)})^{-1/T}
  = \sqrt[T]{\frac{1}{p(\vect x^{(1)}, \ldots, \vect x^{(T)})}} .
\]
\end{nbdef}
To compute it:
\begin{enumerate}[nosep]
  \item compute the probability of the sentence with the tested model and the chain
        rule, $p(\vect x^{(1)}, \ldots, \vect x^{(T)}) = \prod_{t=1}^{T}
        p(\vect x^{(t)} \mid \vect x^{(1)}, \ldots, \vect x^{(t-1)})$;
  \item raise it to the power $-1/T$.
\end{enumerate}
Why both pieces?
\begin{itemize}[nosep]
  \item \textbf{Inverse}: a good model gives the test sentence a \emph{high}
        probability, hence a \emph{low} perplexity. Lower is better.
  \item \textbf{$T$-th root}: the probability of a sentence shrinks with every word, so
        long sentences would always look worse. The $T$-th root is the
        \emph{geometric mean} per word, which makes sentences of different lengths
        comparable.
\end{itemize}

\begin{nbintuition}[perplexity as a number of choices]
\notinslides
A model that spreads its probability uniformly over $k$ words at every step gives the
sentence probability $(1/k)^T$, so its perplexity is
$\bigl((1/k)^T\bigr)^{-1/T} = k$. Perplexity is therefore the \emph{effective number of
equally likely words} the model is hesitating between at each step. A perplexity of 50
means ``as confused as if choosing among 50 words''. A model that knows nothing
(uniform over the whole vocabulary) has perplexity $|V|$.
\end{nbintuition}

\paragraph{Perplexity in practice and its link with the loss.}
\begin{itemize}[nosep]
  \item Computing $\prod_{t=1}^{T} p(\vect x^{(t)} \mid \vect x^{(1)}, \ldots,
        \vect x^{(t-1)})$ for long sequences is \textbf{numerically unstable}: the
        product of thousands of numbers smaller than 1 becomes too small for the
        computer to represent (it rounds to 0). In practice we include begin and end
        markers \texttt{<s>}, \texttt{</s>} and \textbf{truncate the conditioning} at
        sentence or document boundaries, and we work with logarithms.
  \item \textbf{Connection to the cross-entropy loss.} For an RNN, the probability of
        each word is read from the output of the previous step:
        $p(\vect x^{(t+1)} \mid \vect x^{(1)}, \ldots, \vect x^{(t)}) =
        \hat y^{(t)}_{\vect x^{(t+1)}}$. Using $a = e^{\log a}$ and
        $\log \prod = \sum \log$:
        \begin{align*}
          \text{perplexity}
          &= \Bigl[\prod_{t=1}^{T} \hat y^{(t)}_{\vect x^{(t+1)}}\Bigr]^{-1/T}
           = \exp\Bigl( -\frac1T \log \prod_{t=1}^{T} \hat y^{(t)}_{\vect x^{(t+1)}} \Bigr)\\
          &= \exp\Bigl( \frac1T \sum_{t=1}^{T}
             \underbrace{-\log \hat y^{(t)}_{\vect x^{(t+1)}}}_{J^{(t)}} \Bigr)
           = \exp(\text{average cross-entropy loss}) .
        \end{align*}
        Training an LM by minimising the cross-entropy loss is the same as minimising
        perplexity.
  \item \textbf{Lower perplexity} means the model is \textbf{less ``surprised''} by the
        text it is given.
\end{itemize}

\begin{nbwarning}[indices on the slide]
The slide writes $p(\vect x^{(t)} \mid \vect x^{(1)}, \ldots, \vect x^{(t-1)}) =
\hat y^{(t)}_{\vect x^{(t+1)}}$, mixing two conventions, and sums over $i$ while the
terms depend on $t$. The consistent version is the one above: $\hat{\vect y}^{(t)}$ is
computed after reading $\vect x^{(t)}$ and predicts $\vect x^{(t+1)}$, so the $t$-th
factor is $\hat y^{(t)}_{\vect x^{(t+1)}}$. Here the sequence starts with
$\vect x^{(1)} = \texttt{<s>}$ (which is given, not predicted) and the $T$ predicted
tokens are $\vect x^{(2)}, \ldots, \vect x^{(T+1)}$, the last one being \texttt{</s>}.
\end{nbwarning}

\begin{nbexample}{perplexity of two models}
\slidesrc{Language models \& RNNs: 82--84}
The vocabulary is the digits 0 to 9, $|V| = 10$.
\begin{itemize}[nosep]
  \item LM1 outputs any digit with probability $\frac{1}{10}$;
  \item LM2 outputs 0 with probability $\frac{91}{100}$ and each other digit with
        probability $\frac{1}{100}$ (check: $0.91 + 9 \cdot 0.01 = 1$).
\end{itemize}
Which one has the lowest perplexity on the sequence ``0 0 0 3 0'' ($T = 5$)?
\begin{align*}
  \text{LM1:}\quad \text{perplexity} &= \Bigl[\bigl(\tfrac{1}{10}\bigr)^5\Bigr]^{-1/5}
     = \bigl(\tfrac{1}{10}\bigr)^{-1} = 10 \\
  \text{LM2:}\quad \text{perplexity} &= \Bigl[\bigl(\tfrac{91}{100}\bigr)^4 \cdot
     \tfrac{1}{100}\Bigr]^{-1/5} = (0.006857)^{-1/5} \approx 2.709
\end{align*}
LM2 wins: it expects zeros, the sequence is mostly zeros, so it is much less surprised.
LM1 is uniform, so its perplexity is $|V| = 10$, as predicted by the intuition box.
Note that perplexity depends on the test text: on ``3 3 3 3 3'' LM2 would get
$\bigl((1/100)^5\bigr)^{-1/5} = 100$, ten times worse than LM1.
\end{nbexample}

\begin{nbexercise}{2}{perplexity of the toy bigram model}
\notinslides
Compute the perplexity of the bigram model of Section~\ref{sec:ngram-models} (toy
corpus, no smoothing) on the sentence ``I like dogs'', counting the end marker as a
predicted word.
\end{nbexercise}
\begin{nbsolution}
The model predicts $T = 4$ tokens: \emph{I}, \emph{like}, \emph{dogs}, \texttt{</s>},
with probabilities $\frac23$, $1$, $\frac13$, $1$. The sentence probability is
$\frac29$, so
\[
  \text{perplexity} = \bigl(\tfrac29\bigr)^{-1/4} = 4.5^{1/4} \approx 1.46 .
\]
Very low: the toy model has almost no choices to make (after \emph{I} it is certain of
\emph{like}). Lesson: perplexity is only meaningful when comparing models on the
same, realistic, unseen test text; a tiny model evaluated on its own training sentences
looks deceptively good.
\end{nbsolution}

\begin{nbsummary}[RNN language models]
\begin{itemize}[nosep]
  \item $\vect h^{(t)} = \sigma(W_h \vect h^{(t-1)} + W_e E \vect x^{(t)})$,
        $\hat{\vect y}^{(t)} = \softmax(U \vect h^{(t)})$; same weights at every step.
  \item Loss: $J = \frac{1}{T-1}\sum_t -\log \hat y^{(t)}_{\vect x^{(t+1)}}$, minimised
        with SGD on batches.
  \item BPTT: gradient w.r.t.\ a shared weight = sum over its uses,
        $\frac{\partial \vect h^{(t)}}{\partial W_h} \sim \sum_{i} W_h^{\,t-1-i}
        \vect h^{(i)}$; truncated at $\sim$20 steps.
  \item Generation: sample, feed the sample back as input, repeat.
  \item Perplexity $= p(X)^{-1/T} = \exp(\text{average CE loss})$; lower is better;
        uniform model $\Rightarrow$ $|V|$.
\end{itemize}
\end{nbsummary}

\subsection{Problems with RNNs: Vanishing and Exploding Gradients}
\label{sec:vanishing-exploding-gradients}
\slidesrc{Language models \& RNNs: 85--101}

In practice, training the RNN presented so far (often called a \textbf{vanilla RNN},
i.e.\ the plain basic version) is hard, because of two related problems:
\begin{enumerate}[nosep]
  \item \textbf{vanishing gradients};
  \item \textbf{exploding gradients}.
\end{enumerate}
Both come from the same source: in BPTT the gradient travels backwards through many
steps, and at each step it is \emph{multiplied} by the same matrix.

\needspace{12\baselineskip}
\subsubsection{Vanishing gradients: intuition}
\slidesrc{Language models \& RNNs: 87--91}

How much does the loss at step 4 depend on the hidden state at step 1? By the chain
rule, the influence of $\vect h^{(1)}$ on $J^{(4)}$ passes through $\vect h^{(2)}$,
$\vect h^{(3)}$ and $\vect h^{(4)}$, and each passage contributes one factor:
\[
  \frac{\partial J^{(4)}}{\partial \vect h^{(1)}} =
  \frac{\partial \vect h^{(2)}}{\partial \vect h^{(1)}} \times
  \frac{\partial \vect h^{(3)}}{\partial \vect h^{(2)}} \times
  \frac{\partial \vect h^{(4)}}{\partial \vect h^{(3)}} \times
  \frac{\partial J^{(4)}}{\partial \vect h^{(4)}} .
\]

\begin{center}
\begin{tikzpicture}[>=Stealth,
    neu/.style={circle, draw=cyan!70!black, fill=cyan!35, minimum size=5pt, inner sep=0pt},
    hid/.style={draw=cyan!70!black, rounded corners, thick, minimum width=0.5cm,
                minimum height=1.4cm}]
  \foreach \i in {1,...,4} {
    \node[hid] (h\i) at (2.6*\i,0) {};
    \foreach \k in {1,...,4} \node[neu] at (2.6*\i,0.7-0.28*\k) {};
    \node[font=\small, anchor=north] at (h\i.south) {$\vect h^{(\i)}$};
  }
  \foreach \i/\j in {1/2,2/3,3/4} \draw[->, nbgray] ([yshift=4pt]h\i.east) --
      node[above, font=\scriptsize] {$W_h$} ([yshift=4pt]h\j.west);
  \node[font=\small, nbred] (J) at (10.4,1.9) {$J^{(4)}$};
  \draw[->, nbgray] ([xshift=-3pt]h4.north) -- ([xshift=-3pt]J.south);
  % backward arrows getting thinner
  \draw[->, nbred, line width=3pt] ([xshift=4pt]J.south) -- ([xshift=4pt]h4.north);
  \draw[->, nbred, line width=2.2pt] ([yshift=-6pt]h4.west) -- ([yshift=-6pt]h3.east);
  \draw[->, nbred!80, line width=1.3pt] ([yshift=-6pt]h3.west) -- ([yshift=-6pt]h2.east);
  \draw[->, nbred!60, line width=0.5pt] ([yshift=-6pt]h2.west) -- ([yshift=-6pt]h1.east);
  \node[font=\scriptsize, nbred, align=center] at (5.2,-1.45)
       {the gradient shrinks at every step backwards};
\end{tikzpicture}
\captionof{figure}{The gradient of $J^{(4)}$ flows backwards from $\vect h^{(4)}$ to
$\vect h^{(1)}$ (red arrows). If each factor is ``small'', the signal reaching the first
tokens is tiny.}
\end{center}

\begin{nbintuition}[multiplying small numbers]
\notinslides
If each factor were $0.5$, after 3 steps the signal would be $0.5^3 = 0.125$, after 10
steps $0.001$, after 20 steps one millionth. A product of many numbers smaller than 1
\emph{vanishes}; a product of many numbers larger than 1 \emph{explodes}. It is the
same as the children's game ``telephone'': a message whispered along a long chain of
people arrives at the end faint or distorted.
\end{nbintuition}

\begin{nbdef}{Vanishing gradients}
When the factors $\partial \vect h^{(t)} / \partial \vect h^{(t-1)}$ are small, the
gradients with respect to the \textbf{first tokens become tiny}: the loss at a late step
hardly influences how the network processes early words.
\end{nbdef}

\subsubsection{Vanishing gradients: proof sketch}
\slidesrc{Language models \& RNNs: 92--95}

Use again the simplified model $\vect h^{(t)} = W_h \vect h^{(t-1)}$,
$J^{(t)}(\theta) = \mathrm{CE}(\vect h^{(t)}, \vect y^{*(t)})$. Consider the loss at step
$i$ and its gradient with respect to the hidden state of an earlier step $j < i$. Call
$l = i - j$ the \textbf{distance} between the two steps.
\begin{align*}
  \frac{\partial J^{(i)}(\theta)}{\partial \vect h^{(j)}}
  &= \frac{\partial J^{(i)}(\theta)}{\partial \vect h^{(i)}}
     \prod_{j < t \le i} \frac{\partial \vect h^{(t)}}{\partial \vect h^{(t-1)}}
     \reason{chain rule, one factor per step}\\
  &= \frac{\partial J^{(i)}(\theta)}{\partial \vect h^{(i)}}
     \prod_{j < t \le i} W_h
     \reason{$\vect h^{(t)} = W_h \vect h^{(t-1)}$}\\
  &= \frac{\partial J^{(i)}(\theta)}{\partial \vect h^{(i)}}\, W_h^{\,l}
     \reason{$l$ equal factors}
\end{align*}
If $W_h$ is ``small'', the power $W_h^{\,l}$ becomes minuscule as $l$ grows. What
``small'' means for a matrix is made precise by its \textbf{eigenvalues}.

\notinslides
Recall that an eigenvector $\vect q$ of a matrix is a direction that the matrix only
stretches: multiplying by the matrix gives $\lambda \vect q$, where the number
$\lambda$ is the eigenvalue. Multiplying $l$ times stretches it by $\lambda^l$.

Suppose $W_h$ has eigenvalues $\lambda_1, \ldots, \lambda_n$ with
$|\lambda_k| < 1$ and (left) eigenvectors $\vect q_1, \ldots, \vect q_n$ that form a
basis. Write the gradient row vector in this basis,
$\frac{\partial J^{(i)}}{\partial \vect h^{(i)}} = \sum_k c_k \vect q_k\T$. Each
multiplication by $W_h$ multiplies the $k$-th component by $\lambda_k$, so
\[
  \frac{\partial J^{(i)}(\theta)}{\partial \vect h^{(i)}}\, W_h^{\,l}
  = \sum_{k=1}^{n} c_k\, \lambda_k^{\,l}\, \vect q_k\T
  \;\approx\; 0 \quad \text{for large } l ,
\]
because every $\lambda_k^{\,l} \to 0$: the gradient \textbf{vanishes}. How fast?

\begin{center}
\small
\begin{tabular}{@{}c|ccccc@{}}
\toprule
$\lambda$ & $l = 1$ & $l = 5$ & $l = 10$ & $l = 20$ & $l = 50$ \\
\midrule
0.5 & 0.5 & 0.031 & 0.00098 & $9.5 \times 10^{-7}$ & $\approx 0$ \\
0.9 & 0.9 & 0.59  & 0.35    & 0.12                 & 0.0052 \\
1.1 & 1.1 & 1.61  & 2.59    & 6.73                 & 117 \\
\bottomrule
\end{tabular}
\captionof{table}{$\lambda^l$: the factor by which a component of the gradient is
multiplied after travelling back $l$ steps. Below 1 it decays exponentially, above 1 it
grows exponentially.}
\end{center}

The same phenomenon occurs with more realistic RNN models (reference: Pascanu et al.,
\emph{On the difficulty of training recurrent neural networks}, 2013).

\begin{nbwarning}[details the slide leaves out]
\begin{itemize}[nosep]
  \item The condition is on the \emph{absolute value}: $|\lambda_k| < 1$. The slide
        writes $\lambda_k < 1$, but an eigenvalue $-2$ is smaller than 1 and still
        makes the gradient explode. (Eigenvalues can even be complex numbers; then
        $|\lambda|$ is their modulus.)
  \item The argument assumes $W_h$ has a basis of eigenvectors (it is
        \emph{diagonalisable}); the general case is similar, using the largest
        singular value instead.
  \item The slide uses $i$ both for the time step and for the eigenvalue index; above
        the eigenvalue index is $k$.
  \item In a real RNN, $\vect h^{(t)} = \sigma(W_h \vect h^{(t-1)} + \ldots)$, so each
        factor is $\mathrm{diag}\bigl(\sigma'(\cdot)\bigr) W_h$. The derivative of the
        sigmoid is at most $\frac14$ (of $\tanh$ at most 1), which makes the factors
        even smaller and vanishing even more likely.
\end{itemize}
\end{nbwarning}

\subsubsection{Why it matters: long-term dependencies}
\slidesrc{Language models \& RNNs: 96--98}

The gradient that updates the weights is the sum of contributions from all steps. A
contribution from a \textbf{nearby token} arrives almost intact; one from a
\textbf{far-away token} arrives multiplied by $W_h^{\,l}$, i.e.\ almost 0. So the
weights are updated almost only to exploit short-range patterns.

\begin{nbfact}{Long-term dependencies}
With vanishing gradients, \textbf{long-term dependencies are (practically) impossible
to learn}, because the gradient from far-away tokens is much smaller than from nearby
tokens.
\end{nbfact}

\begin{nbexample}{a dependency 20+ tokens back}
Predict the next word:
\begin{quote}\small
After leaving \textbf{work}, he went to pick up his daughter from school. They had
planned to go to the pet shop to buy a hamster and drove through town looking for one
she liked. Eventually they found it, but before heading back home he had to rush back
to his \underline{\quad\quad}
\end{quote}
The right answer is \emph{work}. To predict it the model must connect the blank with
the word \emph{work} at the very beginning, more than 20 tokens earlier. During
training, the gradient that would teach this connection is tiny compared with the
gradients of short-range patterns (``his'' $\to$ a noun, ``rush back to his'' $\to$
\emph{car}, \emph{house}, \ldots), so it provides little signal and learning
long-term dependencies becomes very slow. A vanilla RNN tends to guess from the last
few words only, i.e.\ it behaves much like an n-gram model.
\end{nbexample}

\subsubsection{Exploding gradients and gradient clipping}
\slidesrc{Language models \& RNNs: 99--100}

Symmetrically, if the eigenvalues of $W_h$ are larger than 1 (in absolute value), the
gradients become \textbf{too big}:
\begin{itemize}[nosep]
  \item an update $\theta \leftarrow \theta - \alpha \nabla_\theta J$ with a huge
        gradient is a huge step, which can jump over the minimum and
        \textbf{increase} the loss instead of decreasing it;
  \item in the worst case the numbers overflow and the weights become
        \texttt{inf} or \texttt{NaN} (``not a number''): training is ruined.
\end{itemize}

\begin{nbdef}{Gradient clipping}
If the gradient is too large, scale it down. With a threshold $\tau > 0$:
\begin{enumerate}[nosep]
  \item set $\hat{\vect g} \leftarrow \nabla_\theta J(\theta)$;
  \item \textbf{if} $\norm{\hat{\vect g}} \ge \tau$:
  \item \quad set $\hat{\vect g} \leftarrow \dfrac{\tau}{\norm{\hat{\vect g}}}\,
        \hat{\vect g}$;
  \item \textbf{end if}
\end{enumerate}
then update the weights with $\hat{\vect g}$.
\end{nbdef}

Multiplying by the positive number $\tau / \norm{\hat{\vect g}}$ does not change the
direction of the vector, only its length, which becomes exactly $\tau$. So
\textbf{the update direction remains the same, but the step size is reduced}.

\begin{nbexample}{clipping a gradient}
\notinslides
Gradient $\hat{\vect g} = (30, 40)$, threshold $\tau = 5$. Its length is
$\norm{\hat{\vect g}} = \sqrt{30^2 + 40^2} = 50 \ge 5$, so
\[
  \hat{\vect g} \leftarrow \frac{5}{50}\,(30, 40) = (3, 4),
\]
which has length $\sqrt{9 + 16} = 5 = \tau$ and points in the same direction (the
ratio between the entries is still $3 : 4$).
\end{nbexample}

\begin{center}
\begin{tikzpicture}[>=Stealth]
\begin{axis}[width=8.5cm, height=5.5cm, axis lines=left, xmin=-0.6, xmax=4.6, ymin=0, ymax=5,
    xlabel={parameter $\theta$}, ylabel={loss $J(\theta)$}, xtick=\empty, ytick=\empty,
    label style={font=\small}, clip=false]
  \addplot[domain=-0.3:4.3, samples=200, thick, nbblue]
     {0.15*(x-1)^2 + 3/(1+exp(-8*(x-3)))};
  \fill[black] (axis cs:3,2.1) circle (2pt);
  \draw[->, thick, nbgreen] (axis cs:3,2.1) -- (axis cs:2.5,0.39);
  \node[font=\scriptsize, nbgreen, align=left, anchor=west] at (axis cs:3.15,1.0)
       {with clipping:\\small step, same direction};
  \draw[->, thick, dashed, nbred] (axis cs:3,2.1) .. controls (axis cs:2.0,4.4) ..
       (axis cs:-0.2,4.1);
  \node[font=\scriptsize, nbred, align=center] at (axis cs:1.5,4.7)
       {without clipping: thrown far away};
  \node[font=\scriptsize, nbgray, anchor=west] at (axis cs:3.3,2.3) {cliff};
\end{axis}
\end{tikzpicture}
\captionof{figure}{A ``cliff'' in the loss, typical of RNNs. On the steep wall the
gradient is enormous: a plain gradient step throws the parameters far away, undoing
training; the clipped step goes in the same direction but stays local.}
\end{center}

The slide illustrates this with the 3D loss surface from Pascanu et al.: a flat region
with a steep wall, where unclipped updates bounce off the wall, while clipped ones
proceed smoothly. In practice \textbf{clipping essentially solves the exploding
gradient problem}. It does nothing, however, for vanishing gradients: one cannot
enlarge a gradient that carries no information.

\subsubsection{Solutions for vanishing gradients}
\slidesrc{Language models \& RNNs: 101}

Vanilla RNNs cannot learn long-range dependencies because \textbf{they do not have a
memory}: in
\[
  \vect h^{(t)} = \sigma\bigl(W_h \vect h^{(t-1)} + W_e \vect e^{(t)}\bigr)
\]
the hidden state is \textbf{constantly rewritten}: at every step all of it is
multiplied by $W_h$ and squashed again, so old information fades. The solutions
follow two ideas:
\begin{itemize}
  \item \textbf{Separate memory mechanisms}: \textbf{LSTMs} (Long Short-Term Memory) and
        \textbf{GRUs} (Gated Recurrent Units) add a memory cell and \emph{gates}, small
        learnt switches that decide at each step what to keep, what to forget and what
        to write. Information the gates decide to keep passes from step to step almost
        unchanged, so its gradient does not vanish. These were the
        \textbf{state-of-the-art NLP models in the 2010s}.
  \item \textbf{Direct connections between far-away tokens}: \textbf{attention} lets
        each position look directly at any earlier position, and \textbf{residual
        connections} add a shortcut that skips over layers. The gradient can then
        travel along a short path instead of through dozens of multiplications.
        Attention is the core of the transformer, the architecture of today's LLMs.
\end{itemize}

\begin{nbsummary}[vanishing and exploding gradients]
\begin{itemize}[nosep]
  \item Gradient across $l$ steps: $\frac{\partial J^{(i)}}{\partial \vect h^{(j)}} =
        \frac{\partial J^{(i)}}{\partial \vect h^{(i)}} W_h^{\,l}$ (simplified model).
  \item All $|\lambda_k| < 1$ $\Rightarrow$ $\lambda_k^l \to 0$: \textbf{vanishing}
        gradients, long-term dependencies are not learnt.
  \item Some $|\lambda_k| > 1$ $\Rightarrow$ \textbf{exploding} gradients: loss jumps
        up, \texttt{inf}/\texttt{NaN}.
  \item Exploding $\to$ \textbf{gradient clipping}: if $\norm{\vect g} \ge \tau$, use
        $\frac{\tau}{\norm{\vect g}}\vect g$ (same direction, length $\tau$).
  \item Vanishing $\to$ memory with gates (LSTM, GRU) or direct connections
        (attention, residual connections).
\end{itemize}
\end{nbsummary}

\begin{nbexercise}{3}{vanishing or exploding?}
\notinslides
In the simplified model $h^{(t)} = w\,h^{(t-1)}$ with a scalar hidden state, the
gradient reaching a token $l$ steps back is multiplied by $w^l$. For $w = 0.8$ and
$w = 1.25$, after how many steps does the factor fall below $\frac{1}{100}$ or rise
above $100$? What happens with $w = -1.25$?
\end{nbexercise}
\begin{nbsolution}
For $w = 0.8$ we need $0.8^l < 0.01$, i.e.\ $l > \frac{\ln 0.01}{\ln 0.8} =
\frac{-4.605}{-0.223} \approx 20.6$: from $l = 21$ on (check: $0.8^{21} \approx 0.0092$).

For $w = 1.25$ we need $1.25^l > 100$, i.e.\ $l > \frac{\ln 100}{\ln 1.25} =
\frac{4.605}{0.223} \approx 20.6$: again from $l = 21$ on ($1.25^{21} \approx 108$).
(The two answers coincide because $1.25 = 1/0.8$.)

For $w = -1.25$: $|w^l| = 1.25^l$, the same growth, with a sign that alternates at
every step. The gradient explodes even though $w < 1$, which is why the condition must
be stated on $|w|$ (on $|\lambda|$ for matrices).

Lesson: about 20 steps are enough to lose or blow up the signal by a factor 100, even
with a factor only 20\% away from 1. This is the same order of magnitude as the
truncation length used in BPTT and as the distance in the ``rush back to his work''
example.
\end{nbsolution}

